\chapter{Trigger Finger Release}

\section*{Introduction \& Clinical Indications}
Trigger finger release, also known as stenosing tenosynovitis release, is a common outpatient surgical procedure performed to alleviate the symptoms of trigger finger or trigger thumb. This condition is characterized by the catching, locking, or snapping of a digit when attempting to flex or extend it, often accompanied by pain. The procedure involves surgically dividing the A1 pulley, a fibrous band that forms part of the flexor tendon sheath, typically located at the level of the metacarpophalangeal (MCP) joint in the palm. By enlarging the constricted opening of the tendon sheath, the procedure allows the inflamed and often nodular flexor tendon to glide smoothly without impingement, thereby restoring normal, pain-free digital motion and improving hand function. The surgical intervention is significant as it addresses not only the mechanical obstruction but also the associated pain and functional limitations that can severely impact a patient's quality of life.

\begin{itemize}
  \item Stenosing Tenosynovitis Release
  \item Flexor Tendon Sheath Incision
  \item A1 Pulley Release
  \item Digital Flexor Tenolysis
  \item Percutaneous Trigger Finger Release
  \item Open Trigger Finger Release
\end{itemize}

The underlying principle of trigger finger release is to resolve the mechanical obstruction that prevents smooth gliding of the flexor tendon within its sheath. In stenosing tenosynovitis, the flexor tendon, particularly at the level of the A1 pulley, becomes inflamed and develops a localized nodular thickening. Concurrently, the A1 pulley itself can also thicken, creating a disproportionate size relationship between the tendon and its surrounding tunnel. This mismatch causes the tendon nodule to catch or "lock" as it attempts to pass through the narrowed A1 pulley during finger flexion and extension.

The surgical intervention aims to permanently eliminate this mechanical impediment. By longitudinally incising the A1 pulley, the fibro-osseous tunnel is effectively widened. This enlargement allows the thickened flexor tendon to glide freely and unimpeded, thereby preventing further catching, locking, and associated pain. The technical goals include:
\begin{itemize}
  \item Complete division of the A1 pulley to ensure adequate decompression.
  \item Preservation of the integrity of other crucial pulleys (A2, A3, A4) to prevent tendon bowstringing.
  \item Meticulous protection of adjacent neurovascular structures to avoid iatrogenic injury.
\end{itemize}
Successful release restores the physiological mechanics of tendon excursion, leading to symptom resolution and improved digital function.

The condition now known as trigger finger was first formally described by the French surgeon Alfred Armand Louis Marie Velpeau in 1841. Later, in 1850, Notta provided a more detailed account, attributing the locking phenomenon to a nodule on the flexor tendon. Early treatments were largely conservative, involving massage or splinting. Surgical intervention for trigger finger began to gain traction in the late 19th and early 20th centuries, with surgeons recognizing the mechanical nature of the problem.

The concept of incising the A1 pulley to relieve the constriction became the standard surgical approach. Initial techniques involved open incisions, which have been refined over decades to minimize scar size and improve recovery. In the mid-20th century, the advent of corticosteroid injections provided a less invasive, non-surgical treatment option that proved effective for many patients, often delaying or obviating the need for surgery. More recently, percutaneous release techniques, utilizing a needle or small blade through the skin, emerged as an alternative to open surgery, offering potential benefits of smaller scars and faster recovery, though with specific considerations regarding safety and efficacy. The evolution of treatment reflects a continuous effort to balance efficacy with minimal invasiveness.

Trigger finger release is typically indicated in the following clinical scenarios:
\begin{itemize}
  \item Persistent, painful catching, locking, or snapping of a digit that significantly interferes with daily activities or sleep.
  \item Failure of at least one, and often two, corticosteroid injections to resolve symptoms.
  \item Recurrence of symptoms after initial successful conservative treatment.
  \item Presence of a fixed flexion deformity of the affected digit, indicating chronic tendon impingement.
  \item Palpable, tender nodule at the A1 pulley level with associated triggering.
  \item Multiple digit involvement where conservative measures are impractical or ineffective for all affected digits.
  \item Pediatric trigger thumb or finger, especially if a fixed flexion deformity is present or if spontaneous resolution does not occur by 1 year of age.
  \item Patient preference for definitive surgical treatment after understanding the risks and benefits of all options.
\end{itemize}

Trigger finger release is generally considered a secondary or tertiary treatment option for stenosing tenosynovitis. The primary approach to managing trigger finger typically involves conservative measures, which are less invasive and carry fewer risks. These initial treatments often include rest, activity modification, splinting, non-steroidal anti-inflammatory drugs (NSAIDs), and, most commonly, corticosteroid injections into the flexor tendon sheath. Corticosteroid injections are highly effective for many patients and are usually the first-line medical intervention.

Surgical release is reserved for cases where these conservative treatments have failed to provide sustained relief, or when symptoms are severe, such as a fixed flexion deformity. In some specific circumstances, such as a long-standing, severe fixed flexion deformity, or in certain pediatric cases where spontaneous resolution is unlikely, surgery may be considered earlier. However, for the vast majority of adult patients, it serves as a definitive solution after exhausting non-operative modalities, making it a secondary intervention in the overall treatment algorithm.

\section*{Relevant Surgical Anatomy}
The surgical anatomy relevant to trigger finger release primarily involves the flexor tendons, the A1 pulley, and the surrounding neurovascular structures. The flexor tendons of the fingers run from the forearm through the carpal tunnel and into the fingers, where they are encased in a synovial sheath. The A1 pulley, located at the level of the MCP joint, is a critical anatomical structure that maintains the tendons close to the bone during finger motion. 

The arterial supply to the fingers is provided by the digital arteries, which branch from the superficial palmar arch. Venous drainage occurs through the accompanying veins of the digital arteries. The digital nerves, which provide sensation to the fingers, run adjacent to the flexor tendons and are particularly vulnerable during surgical dissection. 

Key surgical landmarks include:
\begin{itemize}
  \item McBurney's point: Not applicable here, but understanding the MCP joint's location is crucial.
  \item Calot's triangle: Not relevant for this procedure but understanding anatomical relationships is essential.
  \item The A1 pulley: The primary target for surgical intervention.
\end{itemize}

Understanding these anatomical relationships is vital for minimizing complications and ensuring a successful outcome.

\section*{Preoperative Workup \& Preparation}
Preparation for trigger finger release is generally straightforward, especially when performed under local anesthesia. 

\begin{itemize}
  \item \textbf{Medical History and Physical Examination:} A thorough review of the patient's medical history, including allergies, current medications, and any pre-existing conditions (e.g., diabetes, bleeding disorders), is conducted. A focused physical examination of the hand confirms the diagnosis and identifies any associated issues.
  
  \item \textbf{Medication Review:} Patients are typically advised to stop certain medications that can increase bleeding risk, such as anticoagulants (e.g., warfarin, dabigatran) and antiplatelet agents (e.g., aspirin, clopidogrel), for a specified period before surgery. This decision is made in consultation with the prescribing physician and surgeon, weighing the risk of bleeding against the risk of thrombosis.
  
  \item \textbf{NPO Status:} If regional or general anesthesia is planned, patients will be instructed to refrain from eating or drinking for a certain number of hours prior to the procedure (NPO - nil per os). For local anesthesia, this is usually not required.
  
  \item \textbf{Informed Consent:} The surgeon will discuss the procedure in detail, including its benefits, risks, alternatives, and expected recovery, and obtain the patient's informed consent.
  
  \item \textbf{Pre-operative Instructions:} Patients receive instructions on wound care, pain management, and activity restrictions for the post-operative period.
  
  \item \textbf{Antibiotics:} Prophylactic antibiotics are generally not required for routine open trigger finger release due to the low risk of infection, but may be considered in specific high-risk patients or for percutaneous approaches.
\end{itemize}

Routine laboratory tests or imaging studies are typically not necessary unless indicated by the patient's medical history or the type of anesthesia planned.

Careful consideration of contraindications and precautions is essential to ensure patient safety and optimal outcomes.

\textbf{Situations requiring treatment, optimization, or an alternative approach:}
\begin{itemize}
  \item Active infection (e.g., cellulitis, abscess) at the surgical site: Surgery should be delayed until the infection is resolved to prevent spread and complications.
  
  \item Severe, uncontrolled bleeding disorder: The risk of significant hemorrhage outweighs the benefits of elective surgery.
\end{itemize}

\textbf{Relative Contraindications and Precautions:}
\begin{itemize}
  \item Poorly controlled diabetes mellitus: Patients with diabetes have a higher risk of infection, delayed wound healing, and recurrence of trigger finger. Careful glycemic control is advised pre-operatively.
  
  \item Severe peripheral vascular disease: Compromised blood supply can impair wound healing.
  
  \item Severe cardiac or pulmonary disease: May increase the risks associated with any form of anesthesia.
  
  \item Patient inability to cooperate with post-operative care: Essential for proper wound healing and rehabilitation.
  
  \item Proximity of digital neurovascular bundles: Extreme caution is required during dissection, especially in the thumb and little finger, where nerves and vessels are more superficial and prone to injury.
  
  \item Previous surgery in the same area: May lead to altered anatomy and increased scarring, making dissection more challenging.
  
  \item Incomplete release: A risk if the entire A1 pulley is not divided, leading to persistent symptoms.
  
  \item Bowstringing: A rare but serious complication if the A2 or A3 pulleys are inadvertently divided along with the A1 pulley, leading to loss of flexor tendon mechanical advantage.
\end{itemize}

\section*{Surgical Technique \& Procedure}
Trigger finger release is typically performed as an outpatient procedure, often under local anesthesia. The general steps for an open release are as follows:

\begin{itemize}
  \item \textbf{Anesthesia:} The procedure is most commonly performed under local anesthesia, usually a combination of lidocaine with epinephrine injected directly into the palm. This provides excellent pain control and minimizes bleeding. Regional anesthesia (e.g., axillary block) or general anesthesia may be used for patient comfort or in specific circumstances.
  
  \item \textbf{Patient Positioning:} The patient is positioned supine, with the affected hand placed on a specialized hand table. The arm is often prepped and draped to allow for a sterile field.
  
  \item \textbf{Incision:} A small transverse or longitudinal incision, typically 1-2 cm in length, is made in the distal palmar crease directly over the palpable A1 pulley. For trigger thumb, the incision is made at the base of the thumb.
  
  \item \textbf{Dissection and Pulley Identification:} The surgeon carefully dissects through the subcutaneous tissue, taking care to protect the underlying digital nerves and vessels, which lie in close proximity, especially in the thumb and little finger. The A1 pulley, a thickened fibrous band overlying the flexor tendon, is identified.
  
  \item \textbf{Pulley Release:} Using a small scalpel or scissors, the A1 pulley is longitudinally incised from its proximal to its distal attachment. The surgeon ensures complete division of the pulley, often confirming this by palpating the tendon and observing its free glide.
  
  \item \textbf{Confirmation of Release:} The patient is asked to flex and extend the digit to confirm that the tendon now glides smoothly without any catching or locking.
  
  \item \textbf{Hemostasis and Closure:} Any minor bleeding is controlled. The skin incision is then closed with fine sutures or adhesive strips.
  
  \item \textbf{Dressing:} A sterile dressing is applied, often a soft bandage that allows for immediate, gentle finger motion.
\end{itemize}

An alternative, less common technique is percutaneous release, where a needle or small blade is inserted through the skin to divide the pulley without a formal incision. This technique carries a higher risk of incomplete release or neurovascular injury due to its "blind" nature and is generally reserved for specific digits (e.g., middle and ring fingers) with experienced surgeons.

\section*{Complications \& Risks}
While trigger finger release is generally a safe and effective procedure, as with any surgical intervention, there are potential risks and complications. These can be broadly categorized as general surgical risks and procedure-specific risks.

\textbf{General Surgical Risks:}
\begin{itemize}
  \item \textbf{Bleeding:} Minor bleeding is common, but significant hematoma formation requiring re-intervention is rare.
  
  \item \textbf{Infection:} Surgical site infection (SSI) can occur, though it is uncommon for this procedure. Symptoms include increased pain, redness, swelling, and pus discharge.
  
  \item \textbf{Anesthesia Complications:} Risks associated with local, regional, or general anesthesia, including allergic reactions, nausea, vomiting, respiratory issues, or adverse drug reactions.
\end{itemize}

\textbf{Procedure-Specific Risks:}
\begin{itemize}
  \item \textbf{Nerve Injury:} The digital nerves run very close to the A1 pulley. Injury can lead to temporary or permanent numbness, tingling (paresthesia), or pain in the affected digit.
  
  \item \textbf{Vessel Injury:} Damage to the digital arteries can lead to reduced blood flow to the digit, though this is rare.
  
  \item \textbf{Incomplete Release:} If the A1 pulley is not fully divided, the triggering or locking symptoms may persist, potentially requiring a revision surgery.
  
  \item \textbf{Tendon Bowstringing:} A very rare but serious complication if the A2 or A3 pulleys are inadvertently cut. This can lead to a loss of mechanical efficiency of the flexor tendon, resulting in weakness and reduced grip strength.
  
  \item \textbf{Stiffness or Contracture:} Although early motion is encouraged, some patients may experience temporary or, rarely, persistent stiffness of the digit, especially if pre-existing fixed flexion deformity was present.
  
  \item \textbf{Scar Tenderness or Hypersensitivity:} The surgical scar may be tender, itchy, or hypersensitive for several weeks or months.
  
  \item \textbf{Recurrence:} While uncommon after a complete release, symptoms can recur, particularly in patients with underlying systemic conditions like diabetes.
  
  \item \textbf{Complex Regional Pain Syndrome (CRPS):} A rare but severe chronic pain condition that can develop after hand surgery, characterized by disproportionate pain, swelling, stiffness, and skin changes.
  
  \item \textbf{Pillar Pain:} Pain at the base of the palm (the "pillars" of the hand) is a common, usually temporary, complaint after carpal tunnel release, but can occasionally occur after trigger finger release due to surgical manipulation in the palm.
\end{itemize}

\section*{Postoperative Care \& Recovery}
Immediately after trigger finger release, the patient's hand will be bandaged, and they will be monitored in the post-anesthesia care unit (PACU) for a short period. Pain medication will be provided as needed. Most patients are discharged home within a few hours, with instructions for wound care and activity.

During the first 24-48 hours, it is crucial to keep the hand elevated to minimize swelling and to apply ice packs intermittently. Gentle active range of motion (AROM) of the affected finger is encouraged almost immediately to prevent stiffness and promote tendon gliding. Pain is typically managed with over-the-counter analgesics like acetaminophen or ibuprofen, or a short course of prescription pain medication if necessary. The dressing is usually kept clean and dry for the first 24-48 hours.

Over the first 1-2 weeks, the patient will continue with gentle exercises. Sutures, if used, are typically removed at the first follow-up appointment, usually 10-14 days post-operatively. Patients are advised to avoid heavy gripping, lifting, or strenuous activities involving the hand during this period to allow for proper wound healing. Swelling and mild discomfort around the incision site are common and gradually subside. Long-term recovery usually involves a gradual return to full activities over 4-6 weeks, with scar maturation continuing for several months. Most patients experience significant relief of symptoms and restoration of function within this timeframe.

During the surgical procedure, the surgeon documents the condition of the A1 pulley and the flexor tendon. The presence of thickening or nodularity of the tendon, as well as any signs of inflammation, is noted. The pathology report on resected tissues, if applicable, may confirm chronic tenosynovitis or inflammatory changes, although this is uncommon for routine trigger finger release. The primary focus remains on the successful division of the A1 pulley and the restoration of normal tendon function.

A normal, successful postoperative course following trigger finger release is characterized by the resolution of pain and the ability to flex and extend the affected digit without catching or locking. Patients typically report a significant decrease in pain levels within the first few days post-surgery. The surgical site may exhibit mild swelling and tenderness, which usually resolves within a few weeks. Most patients regain full range of motion and return to normal activities within 4 to 6 weeks, with ongoing improvement in function as the surgical site heals.

Post-operative follow-up for trigger finger release is generally straightforward and aims to monitor healing and ensure complete resolution of symptoms. 

\begin{itemize}
  \item \textbf{First Post-operative Visit:} Typically scheduled 10-14 days after surgery. During this visit, the surgeon will:
    \begin{itemize}
      \item Assess the surgical wound for proper healing and signs of infection.
      \item Remove any non-dissolvable sutures.
      \item Evaluate the range of motion of the affected digit and confirm the absence of triggering.
      \item Address any patient concerns regarding pain, swelling, or activity.
    \end{itemize}
  
  \item \textbf{Subsequent Visits:} Usually not required if the recovery is uneventful and the patient is progressing well. However, additional visits may be scheduled if there are concerns about persistent stiffness, pain, or other complications.
  
  \item \textbf{Rehabilitation/Physical Therapy:} For most isolated trigger finger releases, formal hand therapy or physical therapy is not necessary, as gentle active motion is encouraged immediately. However, if significant stiffness persists, or if there was a long-standing fixed flexion deformity pre-operatively, a referral to a hand therapist may be beneficial to guide exercises and improve range of motion.
  
  \item \textbf{Activity Resumption:} Patients are typically advised to gradually increase their activity levels, avoiding heavy gripping or strenuous hand use for several weeks to allow for complete healing and scar maturation.
  
  \item \textbf{Warning Signs:} Patients are instructed to contact their surgeon immediately if they experience any signs of complications, such as increasing pain, redness, swelling, pus discharge from the wound, fever, persistent numbness, or recurrence of triggering.
\end{itemize}

\section*{Outcomes \& Prognosis}
Trigger finger, or stenosing tenosynovitis, is a very common musculoskeletal condition affecting the hand. Its incidence in the general population is estimated to be between 2\% and 3\%. However, this prevalence significantly increases in certain demographic groups and those with specific comorbidities. It is most frequently observed in individuals between 40 and 60 years of age, though it can occur at any age, including in infants (congenital trigger thumb). Women are affected more often than men, with a reported ratio ranging from 2:1 to 6:1. Common indications for trigger finger release often arise from idiopathic causes, but there is a strong association with repetitive gripping activities and underlying systemic conditions. Diabetes mellitus is a significant risk factor, with up to 10\% of diabetic patients experiencing trigger finger, and often presenting with multiple affected digits or recurrence. Other associated conditions include rheumatoid arthritis, gout, and hypothyroidism. The mortality associated with trigger finger release is virtually zero, given its outpatient nature and local anesthesia. Morbidity is generally low, primarily related to the specific surgical risks such as infection, nerve injury, or persistent stiffness, which occur in a small percentage of cases.

Open or percutaneous release usually relieves triggering when the diagnosis is correct, but recovery varies with diabetes, inflammatory disease, duration, stiffness, and prior injections. Persistent pain, stiffness, incomplete release, infection, nerve injury, and recurrence are possible, and some patients require hand therapy or revision. No universal success percentage or recovery time applies.

Recurrence of trigger finger after a complete A1 pulley release is rare, typically occurring in less than 5\% of cases. When recurrence does happen, it is often in patients with underlying systemic conditions such as diabetes mellitus, or if the initial release was incomplete. Prognostic factors influencing outcomes include the patient's overall health, with non-diabetic patients generally having slightly better and faster recoveries. The duration of symptoms and the presence of a fixed flexion deformity pre-operatively can also influence recovery time, with longer-standing deformities sometimes requiring more intensive post-operative therapy to regain full extension. Overall, the prognosis for patients undergoing trigger finger release is excellent, providing a definitive and lasting solution for this common and often debilitating condition.

\section*{References}
\begin{enumerate}
  \item MedlinePlus. Trigger finger release. U.S. National Library of Medicine; 2025.
  
  \item Green's Operative Hand Surgery. 8th ed. Elsevier; 2024.
  
  \item Campbell's Operative Orthopaedics. 14th ed. Elsevier; 2021.
  
  \item American Academy of Orthopaedic Surgeons (AAOS). Clinical Practice Guideline: Management of Trigger Finger. 2023.
\end{enumerate}

\clearpage
